% !TeX root = main.tex
\subsection*{Lighting calculation}
To assess lighting performance, the calculation uses manufacturer LDT intensities given in cd/klm, with each profile scaled by its selected flux: 2843.45 lm at 22 W, 3555.30 lm at 28 W and 4685.55 lm at 38 W. Illuminance is then calculated at each workplane grid cell, with spacing at most 0.2 m, as
\[
E=0.80\sum_j\frac{I_j(C,\gamma)\cos\gamma}{d_j^2}.
\]
A luminaire-to-workplane height of 1.9 m is adopted, corresponding to a 2.7 m ceiling and a 0.8 m workplane. Meanwhile, the maintenance factor of 0.80 assumes periodic cleaning and replacement. Room boundaries use the adopted plan envelopes; each room receives light from its own fittings. Only direct light is included in this calculation.
\begin{center}\begin{tabular}{lrrrr}\toprule
Zone & Lamps & Area m$^2$ & Mean lx & Min/mean\\\midrule
Open office & 67 & 231.87 & 592.9 & 0.136\\
Pantry & 11 & 29.54 & 546.3 & 0.390\\
Upper east room & 10 & 33.68 & 588.5 & 0.267\\
South east room & 4 & 13.07 & 662.6 & 0.175\\
\bottomrule\end{tabular}\end{center}
On this basis, all four mean illuminances exceed the adopted 500 lx target. To provide mounting-height tolerance, the four south-east panels use the 38 W setting. At 2.2 m above the workplane, the respective means are 573.0, 505.8, 545.0 and 589.2 lx; the south-east room retains a 17.8\% margin.

When interpreting these results, the minimum/mean ratios include perimeter cells. Overall, the assessment covers mean direct illuminance; a detailed task-area check would also account for furniture, reflections and glare.
